\chapter{引言}
\label{ch:introduction}

\section{研究背景}

现代游戏智能体需要执行复杂的行为逻辑。对于小型原型，把全部决策写在单一过程式脚本中是可行的；但随着状态和中断条件增加，控制流会越来越难以检查和修改。行为树（Behaviour Tree, BT）把决策逻辑表示为控制节点和叶节点组成的层级结构。复合节点决定如何选择子节点，条件和动作则把抽象决策结构连接到游戏状态与 Actor 代码。高层策略可以独立于移动、动画或战斗方法进行编辑，这是行为树在游戏开发中的主要价值之一\parencite{colledanchise2018bt}。

传统节点--连线树能够直接表达父子拓扑，但会随宽度和深度增长而迅速占用空间。行为树卡片还可能显示参数、描述、Decorator、黑板键和执行状态，比只显示标签的普通层级图更大。当数十或数百张卡片放入有限画布时，开发者必须频繁缩放、平移和搜索，当前执行路径也可能淹没在非活动分支中。这不只是美观问题，因为节点图正是设计者定位任务、理解顺序、修改条件和诊断失败的主要工具。

因此，本项目把一个功能完整的 Godot 行为树插件作为大型树显示研究的平台。在保持上下端口和“同级从左到右执行”的语义基础上，加入可独立开关的焦点+上下文、复杂度管理、导航和运行时解释技术。每项技术都能单独关闭，一方面防止视觉功能故障拖垮编辑器，另一方面保留可复现的基线比较条件。

\section{问题定义与研究空白}

树和图可视化领域已经提出许多方法。鱼眼和焦点+上下文视图会为焦点附近的元素分配更多空间\parencite{furnas1986fisheye,lamping1995hyperbolic}；概览+细节保留全局参考视图\parencite{cockburn2008review}；多列布局可以改善大扇出树的浏览\parencite{song2010largefanout}；增量展开和折叠可以降低图复杂度，同时尽量保持用户的心理地图\parencite{dogrusoz2018complexity,misue1995mentalmap}；大型分层图研究还表明，路径任务的效率取决于具体表示和任务\parencite{bae2011quilts}。

然而，这些研究没有直接回答游戏行为树编辑器应如何组合上述机制。行为树与文件目录等普通层级不同：同级顺序具有执行语义；叶节点会调用代码；Decorator 可能阻止分支运行；图的状态会在每个运行 Tick 中变化。一个可用的方案必须保持执行顺序，支持编辑与撤销，公开黑板状态，并在不向游戏画面加入调试覆盖层的前提下解释失败。现有研究可以指导单项设计，但仍需要一个集成、可执行且可复现评估的 Godot 实现。

\section{研究目标与问题}

本文目标是在一个经过功能验证、能够驱动真实 NPC 的 Godot 4.6 可视化行为树平台上，设计并评估可组合的大型树显示优化。研究问题为：在不同的行为树节点规模和物理屏幕尺寸下，本项目提出的可组合显示优化，相较于传统基线显示，能否改善大型行为树的可测显示与编辑交互表现，同时保持行为树结构和执行语义不变？

\section{研究目标与成功标准}

项目工作分为平台建设与研究评估。平台建设包括：实现可序列化的树、节点、Decorator 和类型化黑板资源；实现 Root、Sequence、Selector、Reactive Selector、Random Selector、Parallel、Repeat、Wait、Action 和 Condition 的 \textsc{success}、\textsc{failure} 与 \textsc{running} 语义；允许可复用组件为 NPC 指定树和 Actor，并由叶节点调用 Actor 方法；提供创建、连接、断开、拖拽、类型化编辑、校验、保存、加载、撤销和重做；只在 Godot 编辑器中显示路径、节点状态、黑板和失败原因。

研究评估为大型树显示方法提供独立持久化开关和安全复位路径，并以原始数据比较不同节点规模、物理屏幕尺寸和优化前后条件下的几何显示、目标定位与上下文保留。在进行显示优化实验前，首先需要确认行为树插件的基础功能正确：各类节点能够按照预期执行，编辑后的行为树能够正确保存和重新打开，撤销与重做也不会破坏树的内容；可玩场景中的复杂敌人应当确实由 241 节点行为树控制。显示优化需要在真实渲染环境中与原始显示方式进行比较。关闭优化后，编辑器应当恢复原来的显示状态，并且不能改变行为树的结构、节点位置和执行顺序。测试过程中只要出现程序报错、崩溃、内存泄漏或异常退出，就认为测试没有通过，即使其他检查项目显示成功也是如此。
